\scanpage{87}
% archive foliation: 86
% lines 1--3
En fait on a ici, exceptionnellement
\[
\widetilde{\unsure{\mathcal T|\Gamma}}_2\simeq\unsure{\mathfrak S_3|\mathcal S_3}\cdot\unsure{\mathcal T|\Gamma}_2
\]
(produit $\frac12$ direct par action de $\unsure{\mathfrak S_3|\mathcal S_3}$ sur $\unsure{\mathcal T|\Gamma}_2$)
% lines 4--6
\diagram{\begin{tikzcd}[ampersand replacement=\&]
\Gamma\arrow[r]\&\deleted{\unsure{\widetilde{\unsure{\mathcal T|\Gamma}}_2}}\ \widetilde{\unsure{\mathcal T|\Gamma}}_2\arrow[r,"\varepsilon",dash]\&\unsure{\mathfrak S_3|\mathcal S_3}
\end{tikzcd}}
\unsure{trivial|triv}

i.e. $\Gamma\longrightarrow\deleted{\unsure{\mathfrak S|\mathcal S}}\widetilde{\unsure{\mathcal T|\Gamma}}_2$, d'où
\[
\boxed{\Gamma\longrightarrow\unsure{\mathcal T|\Gamma}_2^{\unsure{\mathfrak S_3|\mathcal S_3}}}
\]
% lines 7--10
On aurait env. de définir un
iso can. $\unsure{\mathcal T|\Gamma}_2^!\simeq\unsure{\mathcal T|\Gamma}_2^{\unsure{\mathfrak S_3|\mathcal S_3}}$
\margin{?}
et des iso can
\[
\unsure{\mathcal T|\Gamma}_2^!\simeq\unsure{\mathcal T|\Gamma}_\nu^!
\]
% line 11
Pour ceci, \unsure{on} \unsure{va} \unsure{généralison|généralisation} d'
